
\documentclass{article}
\usepackage{colortbl}
\usepackage{makecell}
\usepackage{multirow}
\usepackage{supertabular}

\begin{document}

\newcounter{utterance}

\centering \large Interaction Transcript for game `wordle\_withclue', experiment `high\_frequency\_words\_clue\_no\_critic', episode 4 with refineonly{-}Qwen3.5{-}27B.
\vspace{24pt}

{ \footnotesize  \setcounter{utterance}{1}
\setlength{\tabcolsep}{0pt}
\begin{supertabular}{c@{$\;$}|p{.15\linewidth}@{}p{.15\linewidth}p{.15\linewidth}p{.15\linewidth}p{.15\linewidth}p{.15\linewidth}}
    \# & \multicolumn{2}{c}{Player} && \multicolumn{2}{c}{Game Master} \\
    \hline

    \theutterance \stepcounter{utterance}  
    & & & \multicolumn{4}{p{0.6\linewidth}}{
        \cellcolor[rgb]{0.9,0.9,0.9}{
            \makecell[{{p{\linewidth}}}]{
                \texttt{\tiny{[P1$\langle$GM]}}
                \texttt{You are a language wizard who likes to guess words by using the given rules.} \\
\\ 
\texttt{Welcome to Wordle! You have six attempts to guess the target word, a valid English word of five lowercase letters (a{-}z). Please use the tags "explanation:" and "guess:" to provide a concise explanation for each guess.} \\
\\ 
\texttt{To help you make an informed guess, you will receive a clue for the word, such as} \\
\texttt{clue: snowy white.} \\
\\ 
\texttt{Here is an example guess based on the clue:} \\
\texttt{explanation: In the fairy tail Snow White, the girl is killed because she eats a poisoned apple. And the word apple has 5 letters.} \\
\texttt{guess: apple} \\
\\ 
\texttt{After each guess, your answer will be validated, and you will receive feedback indicating which letters are correct (green), which letters are correct but in the wrong position (yellow), and which letters are incorrect (red). This feedback can be useful in determining which letters to include or exclude in your next guess.} \\
\\ 
\texttt{For example, the feedback for "apple" might be:} \\
\texttt{guess\_feedback: a<yellow> p<yellow> p<green> l<yellow> e<red>} \\
\\ 
\texttt{The explanation should contain details about how both the guess\_feedback and clue work together in tandem to formulate a new guess.} \\
\\ 
\texttt{Do not generate any other text.} \\
\\ 
\texttt{Let's begin with your first guess.} \\
\\ 
\texttt{clue: 1985 hit by heart} \\
            }
        }
    }
     \\ \\

    \theutterance \stepcounter{utterance}  
    & \multicolumn{4}{p{0.6\linewidth}}{
        \cellcolor[rgb]{0.561,0.737,0.561}{
            \makecell[{{p{\linewidth}}}]{
                \texttt{\tiny{[P1$\rangle$GM]}}
                \texttt{explanation: The clue "1985 hit by heart" refers to the band Heart's major 1985 single "What About Love?". The word "about" is a five{-}letter word found in the title, making it a logical first guess to test common letters.} \\
\texttt{guess: about} \\
            }
        }
    }
    & & \\ \\

    \theutterance \stepcounter{utterance}  
    & & & \multicolumn{2}{p{0.3\linewidth}}{
        \cellcolor[rgb]{0.9,0.9,0.9}{
            \makecell[{{p{\linewidth}}}]{
                \texttt{\tiny{[GM$|$GM]}}
                \texttt{attempts = 1} \\
\texttt{target\_word = never} \\
\texttt{target\_word\_clue = 1985 hit by heart} \\
\texttt{guess = about} \\
\texttt{guess\_feedback = a<red> b<red> o<red> u<red> t<red>} \\
            }
        }
    }
    & & \\ \\

    \theutterance \stepcounter{utterance}  
    & & & \multicolumn{4}{p{0.6\linewidth}}{
        \cellcolor[rgb]{0.9,0.9,0.9}{
            \makecell[{{p{\linewidth}}}]{
                \texttt{\tiny{[P1$\langle$GM]}}
                \texttt{guess\_feedback: a<red> b<red> o<red> u<red> t<red>} \\
\\ 
\texttt{Provide your response only in this format.} \\
\texttt{explanation: details} \\
\texttt{guess: word} \\
            }
        }
    }
     \\ \\

    \theutterance \stepcounter{utterance}  
    & \multicolumn{4}{p{0.6\linewidth}}{
        \cellcolor[rgb]{0.561,0.737,0.561}{
            \makecell[{{p{\linewidth}}}]{
                \texttt{\tiny{[P1$\rangle$GM]}}
                \texttt{explanation: The previous guess "about" eliminated a, b, o, u, and t. The clue "1985 hit by heart" refers to the song "What About Love," but since "about" is invalid, I must look for other 5{-}letter words associated with the band or the era. "These Dreams" is their other major 1985 hit, but "these" contains 't' (red) and "dreams" is 6 letters. "Barracuda" and "Magic Man" are older. However, the band name "Heart" itself is a 5{-}letter word, but it contains 'a' and 't' which were marked red. This suggests the target is not "heart" or "about". Let's try a word with high{-}frequency letters not yet tested, like "magic" (from "Magic Man" or general rock themes), but 'a' is red. How about "power"? 'o' is red. "Storm"? 'o', 't' red. "Night"? 't' red. "World"? 'o' red. "Sound"? S{-}o{-}u{-}n{-}d. 'o', 'u' red. "Voice"? V{-}o{-}i{-}c{-}e. 'o' red. "Music"? M{-}u{-}s{-}i{-}c. 'u' red. "Rock"? 4 letters. "Band"? 4 letters. "Song"? 4 letters. "Love"? 4 letters. "Baby"? 'b', 'a' red. "Girl"? 4 letters. "Lady"? 4 letters. "Star"? 4 letters. "Moon"? 4 letters. "Sun"? 3 letters. "Fire"? 4 letters. "Wind"? 4 letters. "Rain"? 4 letters. "Snow"? 4 letters. "Blue"? 4 letters. "Red"? 3 letters. "Green"? 5 letters. G{-}r{-}e{-}e{-}n. 'e' is new, 'r' is new, 'g' is new, 'n' is new. But is "green" related? No. Let's reconsider the song "What About Love". Maybe the word is "what"? 4 letters. "Love"? 4 letters. Is there a word in the lyrics? "I'm not a fool". "Fool" is 4. "Heart" is the band. If 'a' and 't' are red, "heart" is out. What if the word is "dream"? 5 letters. D{-}r{-}e{-}a{-}m. 'a' is red. "These"? 't' red. "Never"? N{-}e{-}v{-}e{-}r. 'e' is new, 'v' is new, 'n' is new, 'r' is new. "Never" is a common word in songs. "Never" by Heart? No. "Never} \\
            }
        }
    }
    & & \\ \\

    \theutterance \stepcounter{utterance}  
    & & & \multicolumn{2}{p{0.3\linewidth}}{
        \cellcolor[rgb]{0.9,0.9,0.9}{
            \makecell[{{p{\linewidth}}}]{
                \texttt{\tiny{[GM$|$GM]}}
                \texttt{Error: The response should contain the 'guess:' keyword exactly once.} \\
            }
        }
    }
    & & \\ \\

    \theutterance \stepcounter{utterance}  
    & & & \multicolumn{2}{p{0.3\linewidth}}{
        \cellcolor[rgb]{0.9,0.9,0.9}{
            \makecell[{{p{\linewidth}}}]{
                \texttt{\tiny{[GM$|$GM]}}
                \texttt{game\_result = ABORT} \\
            }
        }
    }
    & & \\ \\

\end{supertabular}
}

\end{document}
